% =====================================================================
%  Pholama: A 7-Stage Pipeline for Compressing and Specializing a
%           3B-Parameter Large Language Model for On-Device Deployment
%  Publication-grade technical report (compile with pdfLaTeX)
%  Requires: bibtex run for the references (see report.bib placeholder)
% =====================================================================
\documentclass[11pt,a4paper]{article}

% ---------- Packages ----------
\usepackage[utf8]{inputenc}
\usepackage[T1]{fontenc}
\usepackage{lmodern}
\usepackage[margin=1in]{geometry}
\usepackage{amsmath,amssymb}
\usepackage{graphicx}
\usepackage{booktabs}
\usepackage{tabularx}
\usepackage{array}
\usepackage{siunitx}
\usepackage{xcolor}
\usepackage{enumitem}
\usepackage{caption}
\usepackage{subcaption}
\usepackage{float}
\usepackage{fancyhdr}
\usepackage{microtype}
\usepackage[hidelinks,colorlinks=true,linkcolor=blue!50!black,
            citecolor=blue!50!black,urlcolor=blue!50!black]{hyperref}

% ---------- Convenience macros ----------
\newcommand{\param}[1]{\ensuremath{#1\,\text{B}}}
\sisetup{group-separator={,},detect-weight=true,detect-family=true}
\newcolumntype{R}{>{\raggedleft\arraybackslash}X}
\newcolumntype{L}{>{\raggedright\arraybackslash}X}

% ---------- Header / footer ----------
\pagestyle{fancy}
\fancyhf{}
\rhead{Pholama: Compressing \& Specializing a 3B LLM}
\lhead{}
\cfoot{\thepage}
\renewcommand{\headrulewidth}{0.4pt}

% ---------- Title ----------
\title{\textbf{Pholama: A Seven-Stage Pipeline for Compressing and\\
Specializing a 3B-Parameter Large Language Model\\
Toward Ultra-Low-Resource On-Device Deployment}}

\author{
  Nawat Pimolsathien\\
  \small Independent Research Project\\
  \small \texttt{nawatpim@gmail.com}
}
\date{\today}

% =====================================================================
\begin{document}
\maketitle
\thispagestyle{empty}

% =====================================================================
\begin{abstract}
\noindent
Large language models (LLMs) deliver strong natural-language understanding
but their memory and compute footprints preclude deployment on small,
resource-constrained devices. This report presents \textbf{Pholama}, an
end-to-end seven-stage pipeline that compresses and specializes
\emph{Typhoon~2 3B}---a Llama~3.2--based Thai--English instruction model of
\num{3.213} billion parameters---into a domain expert for question answering
about the Office of the Registrar, Chulalongkorn University. The pipeline
combines (i)~\emph{vocabulary trimming} of unused-language tokens, (ii)~%
\emph{structured layer pruning} guided by domain calibration data,
(iii)~\emph{QLoRA recovery fine-tuning}, (iv)~\emph{post-training
quantization} to the GGUF~Q4\_K\_M format, and (v)~\emph{retrieval-augmented
generation (RAG)} at deployment. The resulting model is \textbf{20.5\%
smaller} in parameters (\param{3.21}~$\rightarrow$~\param{2.55}), occupies a
\textbf{1.53~GB} quantized deployment file (a 76.2\% reduction versus the
original bf16 weights), runs \textbf{2.2$\times$ faster} on a bare CPU
(\SI{23.5}{tok/s} vs.\ \SI{10.6}{tok/s} on GPU), and \emph{improves}
in-domain answer quality by 35.7\% relative (\SI{35.0}{\percent}~$\rightarrow$~%
\SI{47.5}{\percent} under RAG). Out-of-domain ability (code/mathematics)
degrades as an \emph{intended} consequence of specialization. The entire
pipeline was executed on a single consumer-grade RTX~3050 Laptop GPU with
\SI{4}{GB} of VRAM, demonstrating that aggressive, research-grade LLM
compression is achievable without cloud infrastructure. We position this work
as a step toward the longer-term objective of running specialized language
models on extremely small devices such as wearables or pocket calculators.
\end{abstract}

\tableofcontents
\bigskip
\hrule
\bigskip

% =====================================================================
\section{Introduction}
\label{sec:intro}

\subsection{Background and Motivation}
The capabilities of large language models scale with parameter count, but so
do their storage, memory-bandwidth, and energy requirements. A typical 3B
model stored in 16-bit precision occupies roughly \SI{6.4}{GB} and requires a
GPU for interactive inference, placing it beyond the reach of phones,
microcontrollers, embedded systems, and other ultra-low-resource hardware.
Yet many practical applications---an institutional FAQ assistant, an offline
help desk, a single-domain chatbot---do not require the model's full general
intelligence. They require fluent comprehension and reasoning \emph{within a
narrow domain}, with factual grounding supplied externally.

This observation motivates a \emph{compress-and-specialize} strategy: rather
than preserving a model's general competence, we deliberately discard the
parts that serve out-of-domain tasks (e.g.\ multilingual coverage, code, and
mathematics) and reinvest the saved capacity into a smaller, faster, domain-%
focused model. The chosen domain is Thai--English question answering for the
Office of the Registrar, Chulalongkorn University (\texttt{reg.chula.ac.th}),
a setting with stable terminology and a well-bounded knowledge base.

\subsection{Objectives}
The project pursues the following concrete objectives:
\begin{enumerate}[label=(O\arabic*),leftmargin=2.4em]
  \item \textbf{Reduce model size.} Shrink the \param{3.21} base model toward
        the \param{1.8}--\param{2.5} range while keeping the deployment file
        within \SIrange{1.5}{2}{GB} of RAM after quantization.
  \item \textbf{Enable on-device / CPU-only inference.} Produce an artifact
        that runs at interactive speed without a discrete GPU.
  \item \textbf{Preserve (and ideally improve) in-domain quality} through
        recovery fine-tuning and retrieval augmentation, while accepting
        out-of-domain regression as evidence of successful specialization.
  \item \textbf{Maintain scientific rigor} via a three-baseline comparison
        framework and reproducible, single-machine execution.
\end{enumerate}

\noindent
The stated \emph{long-term} aspiration of the Pholama project is to push LLM
optimization far enough that a competent, single-purpose language model can
run on extremely small devices---a smartwatch or even a calculator-class
microcontroller. The present report documents the first milestone toward that
vision: a fully functioning, CPU-deployable specialized model derived end-to-%
end on commodity hardware.

\subsection{Contributions}
\begin{itemize}[leftmargin=1.6em]
  \item A reproducible seven-stage compression-and-specialization pipeline,
        each stage executed and validated on a single \SI{4}{GB}-VRAM laptop
        GPU.
  \item An empirical characterization of vocabulary trimming on an
        \emph{English-dominant} tokenizer, correcting the common assumption
        that ``dropping dead languages'' alone yields large savings.
  \item A pruning sensitivity analysis showing that the textbook
        ShortGPT~$28{\to}20$ configuration collapses the model, whereas a
        knee-point of $28{\to}24$ remains recoverable.
  \item A full original-versus-specialized evaluation across three axes
        (size, efficiency, quality) with a deployable RAG demonstrator.
\end{itemize}

% =====================================================================
\section{Methodology}
\label{sec:method}

\subsection{Base Model Specification}
The base model is Typhoon~2 3B, an instruction-tuned Thai--English model built
on the Llama~3.2 3B architecture. Its verified configuration
(Table~\ref{tab:basespec}) was confirmed against the published
\texttt{config.json} prior to any modification.

\begin{table}[H]
  \centering
  \caption{Verified architecture of the base model (Typhoon~2 3B).}
  \label{tab:basespec}
  \begin{tabular}{lr}
    \toprule
    \textbf{Hyperparameter} & \textbf{Value} \\
    \midrule
    Hidden size $d$                       & \num{3072} \\
    Number of transformer layers          & \num{28} \\
    Attention heads                       & \num{24} \\
    Key/value heads (GQA)                 & \num{8} \\
    Head dimension                        & \num{128} \\
    FFN intermediate size                 & \num{8192} \\
    Vocabulary size                       & \num{128256} \\
    Tied input/output embeddings          & \texttt{true} \\
    \midrule
    \textbf{Total parameters}             & \textbf{\num{3212749824}} \\
    \bottomrule
  \end{tabular}
\end{table}

A parameter-budget analysis (Table~\ref{tab:budget}) reveals two facts that
shape the entire compression strategy: (1)~the tied embedding accounts for
only $\sim$12\% of the parameters---so vocabulary trimming helps but is not
the dominant lever; and (2)~within each transformer layer, the feed-forward
network (FFN) dominates at $\sim$75\%, making it the prime target for width
pruning.

\begin{table}[H]
  \centering
  \caption{Where the parameters live in the base model.}
  \label{tab:budget}
  \begin{tabular}{lrr}
    \toprule
    \textbf{Component} & \textbf{Parameters} & \textbf{Share} \\
    \midrule
    Embedding (tied input $+$ output)     & \num{394000000}  & 12.3\% \\
    28 transformer layers                 & \num{2819000000} & 87.7\% \\
    \quad Attention (per layer)           & \num{25200000}   & 25\% of layer \\
    \quad FFN / MLP (per layer)           & \num{75500000}   & 75\% of layer \\
    \bottomrule
  \end{tabular}
\end{table}

\subsection{System Architecture and Pipeline Overview}
The Pholama pipeline (Figure~\ref{fig:pipeline}) transforms the base model
through seven sequential stages. Stages~3--6 monotonically reduce model size;
stage~5 restores quality; stage~7 supplies external factual grounding so that
volatile facts (deadlines, regulations) need not be encoded in the weights.

\begin{figure}[H]
  \centering
  % Replace with an actual diagram, e.g. \includegraphics[width=\textwidth]{figures/pipeline.pdf}
  \fbox{\parbox{0.95\textwidth}{\centering\small
    \texttt{Typhoon2 3B (3.21B)}\\[2pt]
    $\downarrow$ \textbf{P3} vocabulary trim \quad (128{,}256 $\to$ 44{,}926 vocab)\\
    \texttt{vocab-trimmed (2.957B)}\\[2pt]
    $\downarrow$ \textbf{P4} structured layer prune \quad (28 $\to$ 24 layers)\\
    \texttt{pruned (2.554B, ppl 14.1)}\\[2pt]
    $\downarrow$ \textbf{P5} QLoRA recovery fine-tune \quad (ppl 14.1 $\to$ 7.6)\\
    \texttt{recovered (2.554B, bf16)}\\[2pt]
    $\downarrow$ \textbf{P6} quantize $\to$ GGUF Q4\_K\_M \quad (file 1.53 GB)\\
    \texttt{deployable artifact}\\[2pt]
    $\downarrow$ \textbf{P7} RAG (e5-small $+$ FAISS $+$ llama-server)\\
    \texttt{Registrar Q\&A assistant}
  }}
  \caption{The seven-stage Pholama compression-and-specialization pipeline.
  Placeholder for a rendered diagram (\texttt{figures/pipeline.pdf}).}
  \label{fig:pipeline}
\end{figure}

\subsection{Phase 1 --- Environment and Baseline}
Tooling comprises PyTorch (CUDA~12.4), Hugging Face
\texttt{transformers}/\texttt{peft}/\texttt{bitsandbytes}, FAISS, and
\texttt{llama.cpp}. The base model was loaded in 4-bit and profiled to
establish baseline~B1 (latency, VRAM, and parameter count). Parameter counts
are computed from the configuration rather than \texttt{numel()}, because
4-bit packing halves the reported element count.

\subsection{Phase 2 --- Domain Dataset Construction}
A web crawl of \texttt{reg.chula.ac.th} (52 HTML pages and $\sim$30 PDFs)
produced 361 text chunks. A large teacher model (Gemini~2.5 Flash) generated
4{,}198 synthetic Thai--English question--answer pairs grounded in those
chunks, spanning formal/informal registers and Thai/English/code-switched
phrasing. Heuristic filtering (self-reference, no-information, length, and
near-duplicate removal) reduced the set to 4{,}016 clean pairs, split
80/10/10 into \textbf{3{,}212 / 401 / 403} train/validation/test items. A
leak assertion guarantees that no test item appears in calibration or
training data. The same corpus is reused for vocabulary analysis (P3),
pruning calibration (P4), recovery training (P5), and retrieval (P7).

\subsection{Phase 3 --- Vocabulary Trimming}
Following Ushio et al.~\cite{ushio2023vocab}, dead tokens are removed from the
tied \texttt{embed\_tokens}/\texttt{lm\_head} matrix. Running the tokenizer
over the corpus plus train and validation sets (never the test set) yielded
5{,}388 distinct tokens. A key empirical finding contradicts the original
plan's assumption: the Llama~3.2 tokenizer is \emph{English-dominant}
(ASCII/English $\approx$97{,}718 tokens; Thai only 2{,}462; genuinely dead
languages---CJK, Korean, Arabic, Cyrillic---only $\approx$24{,}044). Dropping
dead languages \emph{alone} removes just 18.8\% of the vocabulary. To reach
the target $\sim$40k, an \emph{aggressive} mode additionally drops
rarely-used English/code tokens by tiktoken frequency rank, retaining the
byte alphabet (256), special tokens, all Thai, common Latin-1, punctuation,
corpus tokens, and the lowest-id (most frequent) English tokens, plus the
transitive closure of BPE merges. The result is a vocabulary of
\textbf{44{,}926} tokens (a 65.1\% reduction) and a model of
\textbf{\param{2.9568}}. Any text remains encodable via byte fallback.

\subsection{Phase 4 --- Structured Pruning}
\paragraph{Layer pruning.} Layer importance is estimated with the ShortGPT
Block-Influence criterion---the angular (cosine) distance between a layer's
input and output hidden states under domain calibration data:
\begin{equation}
  \mathrm{BI}(\ell) \;=\; 1 - \mathbb{E}_{x}\!\left[
  \frac{\mathbf{h}^{(\ell)}_{\text{in}}(x)\cdot
        \mathbf{h}^{(\ell)}_{\text{out}}(x)}
       {\lVert \mathbf{h}^{(\ell)}_{\text{in}}(x)\rVert\,
        \lVert \mathbf{h}^{(\ell)}_{\text{out}}(x)\rVert}\right].
  \label{eq:bi}
\end{equation}
Layers with the smallest $\mathrm{BI}$ contribute least and are removed; the
remaining layers are re-indexed (\texttt{layer\_idx} reassignment is required
for correct rotary-position behavior). A sweep over the number of removed
layers (Table~\ref{tab:sweep}) shows a sharp ``knee'': removing 4 layers
($28{\to}24$) yields a recoverable perplexity of 13.2, whereas removing 6 or
8 collapses the model. The plan's textbook $28{\to}20$ target is therefore
\emph{infeasible} for this model.

\paragraph{Width pruning (aggressive variant).} For an even smaller variant,
the FFN intermediate dimension is reduced $8192 \to 6144$ ($-25\%$) using
Wanda-style activation importance, producing a \param{2.1011} model at the
cost of higher perplexity (18.6).

\begin{table}[H]
  \centering
  \caption{Layer-pruning sensitivity sweep (validation perplexity, pre-recovery).}
  \label{tab:sweep}
  \begin{tabular}{lrl}
    \toprule
    \textbf{Layers removed} & \textbf{Resulting depth} & \textbf{Val.\ perplexity} \\
    \midrule
    0 (vocab-trimmed base) & 28 & 5.77 \\
    2  & 26 & 7.2 \\
    \textbf{4} & \textbf{24} & \textbf{13.2 \;(chosen knee)} \\
    6  & 22 & 131 \\
    8  & 20 & 237 \;(output degenerate) \\
    \bottomrule
  \end{tabular}
\end{table}

\subsection{Phase 5 --- QLoRA Recovery Fine-Tuning}
The pruned model is fine-tuned with QLoRA: a 4-bit (nf4) base with LoRA
adapters of rank~16 and $\alpha=32$ applied to the
\texttt{q,k,v,o,gate,up,down} projections. Training used sequence length~512,
batch size~1 with gradient accumulation~16 (effective batch~16), two epochs
(402 steps), a cosine schedule with peak learning rate $2\times10^{-4}$, and
\emph{prompt masking} so that loss is computed only over answer tokens.
Critically, the adapter is \textbf{merged onto a bf16 base}, not the 4-bit
base: merging onto the quantized base corrupts perplexity ($9{\to}17$) and
prevents GGUF conversion. Recovery reduced perplexity from \textbf{14.14 to
7.59} ($-46\%$). Peak training memory was only \SI{2.4}{GB}, well below the
plan's \SIrange{6}{12}{GB} estimate, owing to the small adapter rank,
gradient checkpointing, and a plain AdamW optimizer.

\subsection{Phase 6 --- Quantization}
The recovered bf16 model is converted to GGUF via
\texttt{convert\_hf\_to\_gguf.py} and quantized with \texttt{llama-quantize}
to multiple levels (Q8\_0, Q5\_K\_M, and the deployment target
\textbf{Q4\_K\_M}). The trimmed tokenizer required a one-line patch
(\texttt{res="llama-bpe"}) to satisfy the converter's pre-tokenizer check.
The Q4\_K\_M artifact is \textbf{1.53~GB}, runs at \SI{23.5}{tok/s} on CPU,
and consumes $\approx$\SI{1.9}{GB} of RAM in total---within the target
envelope.

\subsection{Phase 7 --- Deployment and Evaluation}
\paragraph{Deployment (RAG).} The 361 corpus chunks are embedded with
\texttt{intfloat/multilingual-e5-small} (384-dim) and indexed with FAISS. At
query time the top-4 retrieved chunks are injected into the context of a
\texttt{llama-server} instance serving the Q4\_K\_M model
(OpenAI-compatible API). This decouples volatile facts from the weights:
documents can be updated without retraining.
\paragraph{Evaluation.} Quality is measured on a 40-item subset of the test
set, with identical RAG for all conditions, judged by an LLM
(Gemini) using $\text{score}=\text{correct}+0.5\times\text{partial}$.
Efficiency (size, RAM, latency) and out-of-domain regression
(math/code probes) complete the three-axis evaluation.

\subsection{The Three-Baseline Framework}
Per the project's rigor requirement, all results are framed against three
baselines (Table~\ref{tab:baselines}). The present report fully establishes
the \textbf{B1$\leftrightarrow$B3} comparison; the B2 ablation and official
\texttt{lm-eval-harness} runs are identified as remaining work.

\begin{table}[H]
  \centering
  \caption{The three mandatory baselines.}
  \label{tab:baselines}
  \begin{tabularx}{\textwidth}{llL c}
    \toprule
    \textbf{ID} & \textbf{Configuration} & \textbf{Purpose} & \textbf{Status} \\
    \midrule
    B1 & Original Typhoon~2 3B & Upper-bound general model & \checkmark{} measured \\
    B2 & Original $+$ QLoRA $+$ quantize (no prune) & Isolate the prune effect & future work \\
    B3 & Full pipeline (trim $+$ prune $+$ recover $+$ quant.) & The proposed model & \checkmark{} measured \\
    \bottomrule
  \end{tabularx}
\end{table}

\subsection{Reproducibility and Hardware Notes}
All stages ran on Windows~11 with an RTX~3050 Laptop GPU (\SI{4}{GB} VRAM);
bf16 stages were made feasible by CUDA system-memory fallback (spilling to
host RAM, $\sim$9$\times$ faster than CPU). Several platform-specific
pitfalls were encountered and worked around: importing
\texttt{transformers.Trainer}/\texttt{datasets} triggers a PyArrow segfault
(addressed by a hand-written training loop); \texttt{sentence\_transformers}
likewise segfaults (addressed by calling \texttt{AutoModel} directly); and
the GGUF tooling's pip dependencies can silently downgrade PyTorch to a CPU
build.

% =====================================================================
\section{Results and Discussion}
\label{sec:results}

\subsection{Axis 1 --- Size}
Table~\ref{tab:size} summarizes the size reduction. The 20.5\% parameter
reduction decomposes into vocabulary trimming ($-256$M, via dropped
unused-language tokens) and layer pruning ($28{\to}24$, $-4$ layers); the
remaining file-size reduction comes from Q4\_K\_M quantization. Relative to
the original bf16 weights, the deployable artifact is \textbf{76.2\%
smaller}.

\begin{table}[H]
  \centering
  \caption{Axis~1 --- size comparison (original B1 vs.\ specialized B3).}
  \label{tab:size}
  \begin{tabular}{lrrr}
    \toprule
    \textbf{Metric} & \textbf{Original (B1)} & \textbf{New (B3)} & \textbf{Change} \\
    \midrule
    Total parameters             & \num{3212749824} & \num{2554100000} & $-20.5\%$ \\
    Vocabulary                   & \num{128256}     & \num{44926}      & $-65.0\%$ \\
    Transformer layers           & 28               & 24               & $-14.3\%$ \\
    BPE merges                   & \num{280147}     & \num{103362}     & $-63.1\%$ \\
    Weights (bf16)               & \SI{6.43}{GB}    & \SI{5.11}{GB}    & $-20.5\%$ \\
    \textbf{Deploy weights (quant.)} & \SI{1.77}{GB} (4-bit) & \textbf{\SI{1.53}{GB}} (Q4\_K\_M) & $-13.6\%$ \\
    Deploy vs.\ original bf16     & \SI{6.43}{GB}    & \SI{1.53}{GB}    & $\mathbf{-76.2\%}$ \\
    \bottomrule
  \end{tabular}
\end{table}

\subsection{Axis 2 --- Efficiency (Resources and Speed)}
Table~\ref{tab:speed} reports the efficiency results. The headline finding is
that the specialized model on a \emph{bare CPU} (\SI{23.5}{tok/s}) is faster
than the original on a \emph{GPU} (\SI{10.6}{tok/s}), while requiring no
discrete accelerator and $\sim$15.6\% less runtime memory.

\begin{table}[H]
  \centering
  \caption{Axis~2 --- efficiency comparison.}
  \label{tab:speed}
  \begin{tabular}{lrrr}
    \toprule
    \textbf{Metric} & \textbf{Original (B1)} & \textbf{New (B3)} & \textbf{Change} \\
    \midrule
    Runtime RAM            & \SI{2.25}{GB} (VRAM) & $\sim$\SI{1.9}{GB} (RAM) & $-15.6\%$ \\
    Device required        & GPU                  & CPU is enough            & --- \\
    Generation speed       & \SI{10.6}{tok/s} (GPU) & \textbf{\SI{23.5}{tok/s}} (CPU) & $+122\%$ \\
    Prompt speed           & ---                  & \SI{122}{tok/s} (CPU)    & --- \\
    Time / answer ($\sim$60 tok) & $\sim$\SI{6}{s} & $\sim$\SI{3}{s}          & $-50\%$ \\
    \bottomrule
  \end{tabular}
\end{table}

\subsection{Axis 3 --- Quality}
\paragraph{In-domain (improved).} Table~\ref{tab:quality} shows that the
specialized model \emph{outperforms} the original on the domain test set,
both with and without RAG. Recovery fine-tuning taught the model the
domain's answer style and specific facts (e.g.\ Thai request-form codes
such as ``\texttt{Jor.Tor.}'', submission channels, and schedules), whereas
the original answers generically. RAG benefits both models by supplying
ground-truth facts.

\begin{table}[H]
  \centering
  \caption{Axis~3 --- in-domain quality (40-item test subset, LLM judge).}
  \label{tab:quality}
  \begin{tabular}{lrrr}
    \toprule
    \textbf{Condition} & \textbf{Original (B1)} & \textbf{New (B3)} & \textbf{Change (rel.)} \\
    \midrule
    $+$ RAG  & \SI{35.0}{\percent} & \textbf{\SI{47.5}{\percent}} & $+35.7\%$ \\
    No RAG   & \SI{16.2}{\percent} & \SI{41.2}{\percent}          & $+154\%$ \\
    \bottomrule
  \end{tabular}
\end{table}

\paragraph{Out-of-domain (intentionally degraded).} Probes confirm the loss
of general ability (Table~\ref{tab:ood}): arithmetic and code generation fail
after specialization. Per the project's design philosophy, this regression is
not a defect but \emph{positive evidence} that the model traded general
competence for domain expertise and compactness.

\begin{table}[H]
  \centering
  \caption{Axis~3 --- out-of-domain regression probes (new model B3).}
  \label{tab:ood}
  \begin{tabularx}{\textwidth}{lLc}
    \toprule
    \textbf{Task} & \textbf{B3 answer} & \textbf{Status} \\
    \midrule
    $17 \times 23 = ?$       & 411 (correct: 391)                 & wrong \\
    Fibonacci function       & ``uses the recursive method'' (no code) & vague \\
    Solve $2x+5=13$          & $2x=8$ (does not finish $x=4$)      & incomplete \\
    \bottomrule
  \end{tabularx}
\end{table}

\subsection{Discussion}
\paragraph{Where the savings really come from.} The parameter budget
(Table~\ref{tab:budget}) correctly predicted that vocabulary trimming would
be a minor lever for a 3B model ($\sim$12\% of parameters live in the
embedding) and that structured pruning of the transformer stack is the
dominant contributor. The empirical 20.5\% reduction is consistent with this:
vocab trim contributes $\sim$8\% and layer pruning the remainder.

\paragraph{Plan-versus-reality corrections.} Two assumptions in the original
plan did not survive contact with the model. First, vocabulary trimming
cannot reach $\sim$40k by dropping dead languages alone, because the
tokenizer is English-dominant; reaching the target requires deliberately
discarding rarely-used English/code tokens---which is, fortuitously, aligned
with the specialization goal. Second, the canonical ShortGPT $28{\to}20$
configuration is infeasible: removing eight contiguous layers raises
perplexity to 237 and produces degenerate output. The recoverable knee lies
at $28{\to}24$. These corrections are themselves useful contributions, as
both assumptions are common in the compression literature.

\paragraph{Cost-effectiveness.} Combining the three axes, the specialized
model is one-fifth smaller, roughly twice as fast, GPU-free, and
\emph{better} in-domain than the original. For a single-purpose deployment
this is a clearly favorable trade, achieved entirely on commodity hardware.

\paragraph{Threats to validity.} In-domain quality is measured on 40 of 403
test items with an LLM judge---an indicator rather than a certified
benchmark. Out-of-domain regression is, so far, a qualitative probe; official
\texttt{gsm8k}/\texttt{humaneval} numbers are deferred to GPU compute.
Finally, B1 was measured on GPU 4-bit and B3 on CPU Q4, reflecting realistic
deployment conditions but differing hardware. The B2 ablation, which isolates
the pruning contribution from QLoRA$+$quantization, remains outstanding.

% =====================================================================
\section{Conclusion and Future Work}
\label{sec:conclusion}

\subsection{Conclusion}
This report presented \textbf{Pholama}, a complete seven-stage pipeline that
compresses and specializes a \param{3.21} Thai--English LLM into a
\param{2.55}, 1.53~GB, CPU-deployable domain expert for Chula Registrar
question answering. The model is 20.5\% smaller, 122\% faster on CPU than the
original on GPU, fits within $\sim$1.9~GB of RAM, and answers the target
domain 35.7\% better (relative) under RAG---while shedding out-of-domain
ability by design. Every stage was executed on a single \SI{4}{GB}-VRAM
laptop, demonstrating that the full compress-and-specialize workflow is
accessible without cloud resources.

\subsection{Future Work}
\begin{itemize}[leftmargin=1.6em]
  \item \textbf{Complete the baseline matrix.} Run the B2 ablation (original
        $+$ QLoRA $+$ quantize, no pruning) to isolate the pruning
        contribution, and obtain official \texttt{lm-eval-harness}
        (\texttt{gsm8k}, \texttt{humaneval}) numbers across B1/B2/B3 on GPU.
  \item \textbf{Push the aggressive variant.} The \param{2.1011} FFN-pruned
        model is not yet recovered/quantized; with recovery it may offer an
        even smaller deployment point.
  \item \textbf{Toward truly tiny devices.} The project's long-term goal is
        sub-gigabyte, microcontroller-class deployment (smartwatch- or
        calculator-class hardware). Promising directions include extreme
        quantization (2--3 bit, e.g.\ ternary/1.58-bit schemes), KV-cache
        quantization, knowledge distillation into a smaller student, and
        speculative or streaming decoding to reduce latency on constrained
        CPUs.
  \item \textbf{Broaden evaluation.} Expand the human-/judge-scored test set
        beyond 40 items and add calibration and hallucination metrics under
        RAG.
\end{itemize}

% =====================================================================
\section*{Acknowledgments}
\addcontentsline{toc}{section}{Acknowledgments}
The author thanks the maintainers of the open-source tools that made this
work possible---PyTorch, Hugging Face Transformers/PEFT, \texttt{bitsandbytes},
FAISS, and \texttt{llama.cpp}---and acknowledges the public information
published by the Office of the Registrar, Chulalongkorn University, used as
the domain corpus.

% =====================================================================
% References. Two options are provided:
%   (A) BibTeX: uncomment \bibliography lines below and create report.bib
%       (a placeholder report.bib is given in the comment block).
%   (B) A manual thebibliography fallback is included so the document
%       compiles stand-alone without a .bib file.
% =====================================================================
\addcontentsline{toc}{section}{References}

% ---- Option A: BibTeX (recommended) ----
% \bibliographystyle{plain}
% \bibliography{report}
%
% Placeholder report.bib entries:
% @inproceedings{ushio2023vocab,
%   author    = {Ushio, Asahi and Zhou, Yi and Camacho-Collados, Jose},
%   title     = {An Efficient Multilingual Language Model Compression
%                through Vocabulary Trimming},
%   booktitle = {Findings of EMNLP},
%   year      = {2023},
%   eprint    = {2305.15020},
%   archivePrefix = {arXiv}
% }
% @article{men2024shortgpt,
%   author  = {Men, Xin and others},
%   title   = {ShortGPT: Layers in Large Language Models are More Redundant
%              than You Expect},
%   journal = {arXiv preprint arXiv:2403.03853},
%   year    = {2024}
% }
% @inproceedings{dettmers2023qlora,
%   author    = {Dettmers, Tim and Pagnoni, Artidoro and Holtzman, Ari
%                and Zettlemoyer, Luke},
%   title     = {{QLoRA}: Efficient Finetuning of Quantized {LLMs}},
%   booktitle = {NeurIPS},
%   year      = {2023}
% }
% @inproceedings{hu2022lora,
%   author    = {Hu, Edward J. and others},
%   title     = {{LoRA}: Low-Rank Adaptation of Large Language Models},
%   booktitle = {ICLR},
%   year      = {2022}
% }
% @inproceedings{sun2024wanda,
%   author    = {Sun, Mingjie and Liu, Zhuang and Bair, Anna and Kolter, J. Zico},
%   title     = {A Simple and Effective Pruning Approach for Large Language Models},
%   booktitle = {ICLR},
%   year      = {2024}
% }
% @article{dubey2024llama3,
%   author  = {Dubey, Abhimanyu and others},
%   title   = {The Llama 3 Herd of Models},
%   journal = {arXiv preprint arXiv:2407.21783},
%   year    = {2024}
% }
% @article{ma2023llmpruner,
%   author  = {Ma, Xinyin and Fang, Gongfan and Wang, Xinchao},
%   title   = {{LLM-Pruner}: On the Structural Pruning of Large Language Models},
%   journal = {NeurIPS},
%   year    = {2023}
% }
% @misc{lewis2020rag,
%   author = {Lewis, Patrick and others},
%   title  = {Retrieval-Augmented Generation for Knowledge-Intensive NLP Tasks},
%   year   = {2020}, eprint = {2005.11401}, archivePrefix = {arXiv}
% }

% ---- Option B: manual fallback bibliography (compiles without .bib) ----
\begin{thebibliography}{9}

\bibitem{ushio2023vocab}
A.~Ushio, Y.~Zhou, and J.~Camacho-Collados,
``An Efficient Multilingual Language Model Compression through Vocabulary
Trimming,'' in \emph{Findings of EMNLP}, 2023. arXiv:2305.15020.

\bibitem{men2024shortgpt}
X.~Men \emph{et al.},
``ShortGPT: Layers in Large Language Models are More Redundant than You
Expect,'' arXiv:2403.03853, 2024.

\bibitem{dettmers2023qlora}
T.~Dettmers, A.~Pagnoni, A.~Holtzman, and L.~Zettlemoyer,
``QLoRA: Efficient Finetuning of Quantized LLMs,'' in \emph{NeurIPS}, 2023.

\bibitem{hu2022lora}
E.~J.~Hu \emph{et al.},
``LoRA: Low-Rank Adaptation of Large Language Models,'' in \emph{ICLR}, 2022.

\bibitem{sun2024wanda}
M.~Sun, Z.~Liu, A.~Bair, and J.~Z.~Kolter,
``A Simple and Effective Pruning Approach for Large Language Models (Wanda),''
in \emph{ICLR}, 2024.

\bibitem{ma2023llmpruner}
X.~Ma, G.~Fang, and X.~Wang,
``LLM-Pruner: On the Structural Pruning of Large Language Models,''
in \emph{NeurIPS}, 2023.

\bibitem{dubey2024llama3}
A.~Dubey \emph{et al.},
``The Llama 3 Herd of Models,'' arXiv:2407.21783, 2024.

\bibitem{lewis2020rag}
P.~Lewis \emph{et al.},
``Retrieval-Augmented Generation for Knowledge-Intensive NLP Tasks,''
in \emph{NeurIPS}, 2020. arXiv:2005.11401.

\bibitem{llamacpp}
G.~Gerganov \emph{et al.},
``llama.cpp: LLM inference in C/C++,''
\url{https://github.com/ggerganov/llama.cpp}, 2023--.

\end{thebibliography}

\end{document}
